What can we get out of the complementing condition? This question was answered in the work of Agmon, Douglis and Nirenberg \cite{ADN2}. Before we state their results, let us introduce some notation. For $M\in\mathbb{N}$, let
\begin{equation}\label{eq:adn-weighted-sobolev-domain-data-spaces}
\begin{aligned}
X_M&=\prod_{j=1}^kH^{M+t_j}(\Omega),\quad Y_M=\prod_{i=1}^kH^{M-s_i}(\Omega),\\
Z_M&=\prod_{l=1}^mH^{M-r_l-1/2}(\partial\Omega).
\end{aligned}
\end{equation}


We now consider the problem \eqref{eq:weighted-elliptic-general-system} with boundary condition \eqref{eq:weighted-elliptic-general-boundary-conditions}. We write the equations in the compact form $\bm{L}\bm{u}=\bm{f}$ and $\bm{B}\bm{u}=\bm{g}$ and we denote by $\mathcal{A}$ the operator which maps $\bm{u}$ to $(\bm{L}\bm{u},\bm{B}\bm{u}|_{\partial\Omega})$. In choosing weights, we shall now make the convention that $s_i\leq0$ and $t_j\geq0$ for all $i$ and $j$; this can always be achieved by subtracting a constant from the $s_i$ and $r_l$ and adding the same constant to the $t_j$. Let $t\prime=\max_jt_j$, $M_1=\max(0,\max_lr_l+1)$. Then the following holds.

\begin{Theorem}{Agmon, Douglis, Nirenberg}
\label{thm:agmon-douglis-nirenberg-elliptic-estimate}
Let $M\geq M_1$ be an integer. Assume that $\Omega$ is a bounded domain of class $C^{M+t\prime}$, that the coefficients of $L_{ij}$ are of class $C^{M-s_i}(\overline\Omega)$ and that the coefficients of $B_{lj}$ are of class $C^{M-r_l}(\partial\Omega)$. Moreover, assume that ellipticity holds throughout $\overline\Omega$ and that the complementing condition holds everywhere on $\partial\Omega$. Assume that $\bm{f}\in Y_M$ and $\bm{g}\in Z_M$. Then the following hold:\begin{enumerate}\item Every solution $\bm{u}\in X_{M_1}$ is in fact in $X_M$.\item There is a universal constant $K$, independent of $\bm{u}$, $\bm{f}$ and $\bm{g}$, such that, for every solution $\bm{u}\in X_M$, we have
\begin{equation}\label{eq:adn-elliptic-system-apriori-estimate}
\|\bm{u}\|_{X_M}\leq K\left(\|\bm{f}\|_{Y_M}+\|\bm{g}\|_{Z_M}+\sum_{j=1}^k\|u_j\|_{L^2(\Omega)}\right).
\end{equation}
\end{enumerate}If $\bm{u}$ is a unique solution, then the last term in \eqref{eq:adn-elliptic-system-apriori-estimate} can be omitted.
\end{Theorem}

The result thus consists of a regularity statement and an a priori estimate. Agmon, Douglis and Nirenberg actually prove more than we have stated; they establish similar results in $L^p$-based Sobolev spaces and also in Hölder spaces. We also note that some of the smoothness hypotheses on $\Omega$ and the coefficients can be weakened. We shall not pursue this point here. A proof of the theorem is beyond the scope of this introductory text. However, we refer to Sections~\ref{sec:elliptic-interior-regularity} and \ref{sec:elliptic-boundary-regularity} for a proof of a special case, namely, second-order elliptic PDEs with Dirichlet boundary condition. We next derive an interesting corollary.

\begin{Corollary}
\label{cor:adn-boundary-operator-semi-fredholm}
Let all assumptions be as in the preceding theorem. Assume in addition that $M+t_j>0$ for every $j$. Then the operator $\mathcal{A}:X_M\to Y_M\times Z_M$ is semi-Fredholm.
\end{Corollary}

\begin{Proof}
It easily follows from the smoothness hypotheses on the coefficients that $\mathcal{A}$ does indeed map $X_M$ to $Y_M\times Z_M$. Let $N(\mathcal{A})$ be the nullspace of $\mathcal{A}$, and let $B$ be the intersection of $N(\mathcal{A})$ with the unit ball in $(L^2(\Omega))^k$. By the theorem, $B$ is bounded in the norm of $X_M$, hence precompact in $(L^2(\Omega))^k$. Since the unit ball in an infinite-dimensional space is never precompact, $N(\mathcal{A})$ must be finite dimensional.

Next, we shall show that the range of $\mathcal{A}$ is closed. For that purpose, assume that $\bm{u}_N$ is a solution of $\bm{L}\bm{u}_N=\bm{f}_N$ with boundary conditions $\bm{B}\bm{u}_N=\bm{g}_N$, and that $\bm{f}_N$ and $\bm{g}_N$ converge in $Y_M$ and $Z_M$ to $\bm{f}$ and $\bm{g}$, respectively. Without loss of generality, we may assume that $\bm{u}_N$ is perpendicular to $N(\mathcal{A})$ in $(L^2(\Omega))^k$. We claim that $\bm{u}_N$ is then bounded in $(L^2(\Omega))^k$. Suppose not. After taking a subsequence, we may assume $\|\bm{u}_N\|_2\to\infty$. Let $\bm{v}_N=\bm{u}_N/\|\bm{u}_N\|_2$. Then $\bm{v}_N$ solves the problem $\bm{L}\bm{v}_N=\bm{f}_N/\|\bm{u}_N\|_2$ with boundary conditions $\bm{B}\bm{v}_N=\bm{g}_N/\|\bm{u}_N\|_2$. It follows from \eqref{eq:adn-elliptic-system-apriori-estimate} that the sequence $\bm{v}_N$ is bounded in $X_M$. Hence it has a subsequence which converges weakly in $X_M$, hence strongly in $(L^2(\Omega))^k$. Let $\bm{v}$ be the limit. Then $\bm{v}$ is in the nullspace of $\mathcal{A}$ and in its orthogonal complement, hence zero. But this is a contradiction, since $\|\bm{v}\|_2=\lim_{n\to\infty}\|\bm{v}_N\|_2=1$. Since $\bm{u}_N$ is bounded in $(L^2(\Omega))^k$, \eqref{eq:adn-elliptic-system-apriori-estimate} implies that it is also bounded in $X_M$. Hence after taking a subsequence, $\bm{u}_N$ converges weakly in $X_M$ and strongly in $(L^2(\Omega))^k$. Applying \eqref{eq:adn-elliptic-system-apriori-estimate} again, we see that $\bm{u}_N$ actually converges strongly in $X_M$. The limit $\bm{u}$ is a solution of $\bm{L}\bm{u}=\bm{f}$ with boundary conditions $\bm{B}\bm{u}=\bm{g}$.
\end{Proof}
